\section{Introduction}
\label{sec:intro}

Geysering is the rapid expulsion of an air--water mixture through a manhole,
dropshaft, or ventilation shaft during transient filling of a drainage
system. The resulting jet can displace covers, damage nearby infrastructure,
and endanger the public. Field pressure records rule out a purely inertial
water-column explanation because the piezometric head near an erupting shaft
need not reach ground level \cite{wright2011note}. Controlled experiments
instead identify the release and expansion of a large entrapped air pocket
through a partly or fully water-filled shaft as the principal driving
mechanism \cite{vasconcelos2011,wright2011note}.

Laboratory evidence now spans several manifestations of this mechanism.
Vasconcelos and Wright \cite{vasconcelos2011} isolated air-pocket release in
a single ventilation shaft and documented both benign venting and geysering.
Cong et al.\ \cite{cong2017} varied riser diameter and initial pocket volume
systematically and proposed a two-parameter regime criterion. Later studies
examined larger shafts \cite{muller2017}, prototype-motivated junction
chambers \cite{liu2020}, covered manholes \cite{zhang2022}, and mitigation
measures \cite{qian2020}. Related multiphase-flow work has also isolated the
rapid expansion of a Taylor bubble in a vertical liquid column
\cite{zhouprosperetti2019}. Together, these studies show that geysering is not
set by a pressure threshold alone. It emerges from competition among pocket
compression, gas delivery to the shaft, counter-current exchange, and the
time available for the water column to reach the rim.

Available numerical approaches resolve different parts of this competition.
Three-dimensional volume-of-fluid simulations can reproduce the local
interface and jet structure of an individual event \cite{chan2018}. Their
cost, however, limits their use for network-scale studies involving many
shafts and rainfall scenarios. One-dimensional mixed-flow solvers are more
suited to that scale, but their gas representation is often prescribed.
Preissmann-slot and two-component-pressure models primarily follow liquid
pressurization \cite{vasconcelos2006}. Entrapped gas is then introduced
through a lumped pocket law, a specified release rate, or a boundary
condition. Rigid-column two-phase models likewise require the gas-pocket and
water-column topology to be known in advance \cite{vasconcelos2011}. Such
models can interpret a selected apparatus, but they cannot independently test
whether the computed gas--liquid dynamics select venting or eruption.

A companion methods paper \cite{xue2026model} addressed this modelling gap
for transient mixed flow in closed conduits. Its stratified branch conserves
gas mass and gas momentum within a decoupled two-fluid system. Its liquid-full
branch uses an elastic-pipe closure, while moving pressurized--stratified
interfaces are solved as Regime-Transition Riemann Problems. A conservative
cut-cell update tracks these fronts without prescribing their motion. The
methods paper established the numerical formulation through canonical
mixed-flow tests and included a schematic geyser-onset example. It did not
quantitatively assess the framework against geysering experiments.

The present study provides that application-level assessment. The companion
framework is extended to each experimental layout by representing the shaft
as the same one-dimensional two-fluid system rotated into the vertical. A
local junction closure connects the shaft to the horizontal conduit. The
closure is conditioned on the resolved flow regime and geometry, but its
parameters and the numerical configuration are frozen before comparison.
Consequently, gas release, water-column displacement, and regime selection
are computed outcomes. No measured release history or geyser trajectory is
imposed.

The assessment is organized around three complementary evidence sets. The
first campaign compares time-resolved pressure, interface, and free-surface
histories for the non-geysering and geysering ventilation-shaft tests of
Vasconcelos and Wright \cite{vasconcelos2011}. The second treats the model as
a regime classifier against the seven filmed riser tests and the
63-configuration parameter sweep of Cong et al.\ \cite{cong2017}. The third
examines the junction-chamber experiments of Liu et al.\ \cite{liu2020},
including a downstream-boundary branch reversal and a transient mass-
oscillation event. These campaigns test increasingly complex questions:
whether the model selects the correct branch, whether that selection
generalizes across a regime map, and whether the reduced representation
retains the first-order dynamics of a prototype-motivated chamber.

This paper therefore asks a narrower question than the companion methods
study: how far can a gas-conserving one-dimensional framework reproduce
observed geysering before unresolved junction, film, and acoustic physics
become controlling? Section~\ref{sec:model} summarizes the formulation and
defines the application-specific closures. Section~\ref{sec:tests} reports
the three experimental comparisons. Section~\ref{sec:discussion} interprets
the common successes and failure modes, and Section~\ref{sec:conclusions}
states the resulting scope of use.
